% ---------------------------------------------------------------------
% 2.1 Contexto: discapacidad motora severa y comunicación
% ---------------------------------------------------------------------
\section[DISCAPACIDAD MOTORA SEVERA Y COMUNICACIÓN]{Discapacidad motora severa y comunicación}

De acuerdo con el Informe mundial sobre la discapacidad elaborado por la Organización Mundial de
la Salud en conjunto con el Banco Mundial, se estima que más de mil millones de personas en el
mundo viven con alguna forma de discapacidad, de las cuales cerca de 200 millones presentan
dificultades considerables para realizar actividades que se consideran básicas en la vida diaria
\cite{oms2011}. En el caso de México, la Encuesta Nacional de la Dinámica Demográfica de 2018
reportó una prevalencia de discapacidad del 6.3\% de la población, es decir, alrededor de 7.8
millones de habitantes, de los cuales el 17.8\% tenía dificultad para mover o usar brazos o manos y
el 10.5\% presentaba dificultad severa para hablar o comunicarse \cite{inegi2018}. Aunado a lo
anterior, Feigin et al. señalan que los trastornos neurológicos son actualmente la principal causa
de discapacidad a nivel mundial y que su carga tiende a crecer con el envejecimiento de la
población \cite{feigin2020}.

Dentro de este grupo se encuentran padecimientos como la esclerosis lateral amiotrófica (ELA), las
lesiones medulares altas o los eventos vasculares del tallo cerebral, que en sus etapas avanzadas
pueden llevar al llamado síndrome de enclaustramiento (\textit{locked-in}), condición en la que la
persona conserva sus funciones cognitivas pero pierde prácticamente todo el control muscular
voluntario, incluido el del habla. Uno de los primeros dispositivos de deletreo controlados
únicamente con la actividad cerebral fue presentado por Birbaumer et al. precisamente para
pacientes con ELA en esta condición, quienes lograron escribir mensajes a partir del control
voluntario de potenciales corticales lentos, aunque tras semanas de entrenamiento y a una
velocidad de escritura considerablemente baja \cite{birbaumer1999}.

La comunicación aumentativa y alternativa agrupa las estrategias y herramientas que complementan o
sustituyen el habla, desde tableros de letras hasta sistemas de seguimiento ocular. Sin embargo,
casi todas ellas dependen de algún canal muscular residual, ya sea el movimiento de los ojos, de
un dedo o de los párpados, por lo que dejan de ser útiles cuando dicho control se pierde por
completo. Es en este punto donde las interfaces cerebro-computadora (BCI) ofrecen un canal de
salida que no depende de los nervios periféricos ni de los músculos \cite{wolpaw2002}, y en este
sentido Sellers y Donchin realizaron las primeras pruebas de un sistema basado en la componente
P300 con pacientes con ELA, mostrando que estos podían operarlo a pesar del avance de la
enfermedad \cite{sellers2006}.
